% Front matter only. The title, authors and affiliation live in main.tex, where
% LNCS expects them -- putting them here printed the author names into the body of
% the paper, on top of the anonymous block, which breaks double-blind review.

\begin{abstract}
A web application firewall must commit, allow or block, on the evidence a single
request carries, and deception normally happens \emph{after} that commitment: a
session already judged hostile is moved into a honeypot. We treat deception instead
as a move available while the detector is still undecided. When intent is uncertain,
the defence adds an invisible, inert probe to the response: a fake table named in an
error, an unused JSON field, a hint at a deprecated endpoint. An honest client never
renders it; a probing one acts on it. The contribution is the rule that decides when
to deploy one. Probing has no immediate benefit, since the request still reaches the
real application, so its worth is entirely the information a bite reveals. Pricing
that worth as the expected value of sample information makes probing cost-optimal
over a belief band whose edges are outputs of a frozen cost table and a measured
bite likelihood ratio; under cost accounting alone the band is empty, so there is no
middle action to tune. Over 100 seeded traffic draws against one frozen model, a
randomised holdout estimates the probe's causal effect at $+0.051$
$[+0.034, +0.068]$ (Fisher exact, $p < 10^{-5}$), and recall rises from 0.917 to
0.951 (paired McNemar, $p < 10^{-4}$), concentrated in the one subcategory where the
passive detector is undecided. None of 7{,}920 benign sessions is diverted.

\keywords{Cyber deception \and Honeypots \and Value of information \and
Cost-sensitive detection \and Web application security \and Honeytokens}
\end{abstract}
